\documentclass[11pt]{article}
\usepackage[margin=1in]{geometry}
\usepackage[T1]{fontenc}
\usepackage{lmodern}
\usepackage{microtype}
\usepackage{amsmath,amssymb}
\usepackage{booktabs}
\usepackage{tabularx}
\usepackage{array}
\usepackage[numbers,sort&compress]{natbib}
\usepackage{xcolor}
\usepackage[colorlinks=true,linkcolor=blue!60!black,citecolor=blue!60!black,urlcolor=blue!60!black]{hyperref}
\usepackage{enumitem}
\setlist{nosep,leftmargin=*}
\newcommand{\anibench}{\textsc{anibench}}
\newcommand{\y}{\checkmark}
\newcommand{\n}{--}
\newcolumntype{C}{>{\centering\arraybackslash}X}

\title{Positioning \anibench{} in the 2026 Literature:\\[2pt]
\large A Second-Pass Review and Decision Memo}
\author{Draft for the author team \quad (prepared 17 November 2026)}
\date{Literature snapshot as of 17 November 2026. Internal; not for distribution.}

\begin{document}
\maketitle

\begin{abstract}
We conducted a second, more exhaustive review of work related to \anibench{}, covering code-based animation, regional image animation, localized editing, agent side effects, verifiable visual rewards, part grounding, and grader validity, with emphasis on papers released in 2025 and 2026. The review adds 58 verified references to the 58 of the first pass. Its central finding is that both fields from which \anibench{} draws its claims filled substantially during 2026. Deterministic verification of code-generated animation is now established, most directly by MotionGraphicsBench, released ten days before this memo~\citep{ma2026motiongraphicsbench}, and deterministic scoring of side effects now exists for agent actions, SVG, web pages and code~\citep{huang2026sabench,pysklo2026agentdiff,gupta2026vectorbench,li2026mtweb2code,zhu2026overediting}. Consequently, ``judge-free'' and ``separates change from side effects'' no longer distinguish \anibench{} on their own. We find no prior work, however, that measures the side effects of a \emph{motion} applied to a \emph{raster} image, in which the model must infer the part, its joint, and the region it may sweep, and in which a side effect may be a few pixels or the part coming apart. We recommend retaining the focused benchmark and making its headline the phenomena that model judges and preference arenas cannot detect, an analysis that only a grader validated against sized corruptions permits. We list three claims the paper should no longer make.
\end{abstract}

\section{Method of This Review}
A research agent searched arXiv, proceedings, project pages, lab release notes, leaderboards, and practitioner forums between 17 and 18 November 2026. Every cited paper was opened and its title, authors, date and venue were copied from the opened page; venues are given only where the page states one. The authors independently re-verified the two papers most consequential for positioning (MotionGraphicsBench and SABench) against their arXiv pages. Two limitations apply. First, posts on X (formerly Twitter) and Reddit could not be accessed programmatically, so practitioner sentiment rests on secondary sources: news coverage, newsletters, and Hacker News. Second, scoring details for some benchmarks (Frontier3DBench, DLEBench) were not stated on the pages opened, and we do not characterize them beyond what was stated.

\section{The Field as of November 2026}

\paragraph{Deterministic verification of code-generated animation.} MotionGraphicsBench gives a model a static SVG that already contains every element, together with a motion prompt, and asks it to animate the SVG; it comprises 100 instances derived from real Lottie designs~\citep{ma2026motiongraphicsbench}. Each instance is graded by a deterministic verification program that accepts different valid implementations. On a balanced set of 180 animations, the programs agree with human judgments 93.3\% of the time, against 83.3\% for a VLM judge. Thirty-two LLMs in 47 configurations achieve pass@1 between 6.0\% and 67.7\%. Its predecessor, MoVer, verifies spatio-temporal properties of SVG animations in first-order logic~\citep{ma2025mover}. Neither scores side effects as a separate, always-on quantity: according to the agent's reading of the full paper, preservation enters MotionGraphicsBench only as per-requirement checks on some prompts. Other 2026 benchmarks generate animation from scratch and score it against a reference or with judges~\citep{zhang2026animwebbench,diao2026vibevideobench,wang2026simuscene}.

\paragraph{Deterministic side-effect scoring.} In agent evaluation, at least nine benchmarks released during 2026 score unrequested or out-of-scope actions with deterministic oracles. SABench, accepted to the NeurIPS 2026 Evaluations and Datasets Track, finds that 15.2\% of 2{,}211 model-task runs incur collateral damage and that 12.8\% of runs satisfying the goal still perform unauthorized writes~\citep{huang2026sabench}; Agent-Diff fails any run whose final state differs from the goal state outside the requested change~\citep{pysklo2026agentdiff}. In artifacts, Vector-Bench requires SVG repairs to land while 61 protected objects per task remain unchanged, and the best of 34 systems reaches 15.0\% full success~\citep{gupta2026vectorbench}; MT-Web2Code scores target fidelity and preservation of unaffected web content on separate axes~\citep{li2026mtweb2code}; and an over-editing benchmark measures excess edits in code against a known minimal patch~\citep{zhu2026overediting}.

\paragraph{Preservation in regional image-to-video generation.} Methods that animate a selected region measure static-region preservation with user studies, whole-frame consistency, or tracker distances~\citep{ma2024followyourclick,singer2025timetomove}. SNF-Bench separates static-region fidelity from motion in fixed-camera video, validates each metric by injecting known corruptions, and shows that whole-frame motion metrics increase by only 1.07$\times$ under maximal injected drift while its static-region metrics increase by 1.32 to 1.86$\times$~\citep{minar2026snfbench}. This is the argument \anibench{} makes against whole-frame similarity, reached independently in a different field with the same validation logic.

\paragraph{Grader validity.} VerifierBench finds that LLM-written verifiers reject invalid solutions more reliably than they accept valid alternatives, and recommends testing verifiers against other admissible solutions~\citep{yin2026verifierbench}; our study of 242 alternative valid animations is an instance of that recommendation. A construct-validity review of 445 benchmarks asks authors to state what claim their scores support~\citep{bean2025constructvalidity}, a criterion that ICML 2027 now applies to its Datasets, Evaluations and Benchmarks contribution type. Two-sided grader validation is therefore becoming an expectation rather than a differentiator.

\paragraph{Part grounding.} Multimodal models ground whole objects well but degrade when the query names a part~\citep{mehrab2026ophrg}, and pointing benchmarks measure localization directly~\citep{cheng2025pointarena}. These results support the hypothesis that \anibench{} failures originate in grounding, and suggest a diagnostic probe described in Section~\ref{sec:rec}.

\begin{table}[t]
\centering
\small
\setlength{\tabcolsep}{3pt}
\begin{tabularx}{\linewidth}{lllCCCC}
\toprule
Benchmark & Input & Output & Change verified & Side effects scored separately & Integrity / pivot & Grader validated against \\
\midrule
MotionGraphicsBench & SVG & GSAP & \y & partial & \n & humans (93.3\%) \\
MoVer & SVG & SVG anim. & \y & \n & \n & synthetic prompts \\
Vector-Bench & SVG & SVG & \y & \y\ (static) & \n & \n \\
MT-Web2Code & web page & HTML/CSS & \y & \y\ (static) & \n & \n \\
SABench & tool state & actions & \y & \y\ (discrete) & \n & oracle design \\
SNF-Bench & video & video & \y & \y\ (frames) & \n & injected corruptions \\
\midrule
\anibench{} & raster & three.js & \y & \y\ (every pixel, every frame) & \y & corruptions, valid alternatives, humans (planned) \\
\bottomrule
\end{tabularx}
\caption{Closest neighbors, by the properties that define \anibench{}. ``Partial'' indicates that preservation is checked only as a requirement of some prompts.}
\label{tab:neighbors}
\end{table}

\section{What Remains Distinctive}
Table~\ref{tab:neighbors} summarizes the comparison. We found no benchmark in which a model animates a raster image by writing code, and no image-to-video evaluation with a pixel-exact static-region test. The distinctive combination is: (i) a raster input, in which the model must infer the part, its joint and its swept region, whereas in SVG the element boundaries are given; (ii) preservation scored on every pixel of every frame against a region larger than the part; (iii) part integrity, connectivity and pivot checks, which have no counterpart in the works above; and (iv) a single grader with fixed thresholds validated against seeded corruptions of known size and against independent valid solutions. We propose the following positioning statement:

\begin{quote}
Benchmarks now verify requested changes deterministically and score side effects deterministically in discrete state. None measures the side effects of a motion on a raster image, where the model must infer the part, its joint and the region it may sweep, and where a side effect may be a few pixels or a part coming apart. \anibench{} measures this, separately from whether the motion occurred, with a grader validated in both directions.
\end{quote}

\paragraph{Claims to retract.} The following statements in our current drafts are contradicted by 2026 work and should be removed: ``no benchmark separates change from preservation''~\citep{gupta2026vectorbench,li2026mtweb2code,huang2026sabench,minar2026snfbench}; ``animation benchmarks rely on judges or similarity''~\citep{ma2026motiongraphicsbench,ma2025mover}; and ``side effects of agents are unmeasured''~\citep{pysklo2026agentdiff,huang2026sabench}.

\section{Candidate Directions}
\label{sec:rec}

We evaluated six directions for the nine weeks before the ICML deadline (Table~\ref{tab:dir}). The assessments are judgments from the evidence above, not measurements.

\begin{table}[t]
\centering
\small
\setlength{\tabcolsep}{3pt}
\begin{tabularx}{\linewidth}{>{\raggedright\arraybackslash}p{4.2cm}lll>{\raggedright\arraybackslash}X}
\toprule
Direction & Novelty & Fit to Jan 22 & Risk & Assessment \\
\midrule
(a) Focused benchmark, polished & medium & good & medium & Unique intersection, but reviewers will compare with MotionGraphicsBench first. \\
(b) Side-effect-free edits across modalities & low per modality & poor & high & Each modality now has a 2026 benchmark; defer to a later venue. \\
(c) RL environment with training gains & medium-high & poor to fair & high & No visual work rewards preservation, but rewards are sparse for weak models; reduce to a reward-validity claim. \\
(d) Image-to-video comparison & medium-high & fair & medium & Requires a calibrated noise floor, since video models re-encode every pixel. \\
(e) Agent-safety framing & low & good & high & Crowded; our failures reflect localization capability, not policy. \\
(f) What judges and arenas miss & high & good & low to medium & Uses sized corruptions as free ground truth; differentiates from MotionGraphicsBench with a result. \\
\bottomrule
\end{tabularx}
\caption{Candidate directions.}
\label{tab:dir}
\end{table}

Direction (f) consists of three experiments built from assets already in the plan. \emph{Judge blindness:} because each seeded corruption has a known size (a 20\,px fragment, a 2\,px seam, a moved background), we can present corrupted and valid animations to frontier VLM judges and plot detection rate against side-effect size. Prior findings that judges award credit for features absent from the running program~\citep{bakshi2026worldbench} predict that judges miss small side effects that the pixel checks detect. No other benchmark has labeled corruptions at controlled sizes, so this analysis is unavailable to competing work. \emph{Arena divergence:} the rank correlation between \anibench{} and public preference arenas (planned experiment E8) tests whether those arenas measure locality at all. \emph{Failure attribution:} combining the procedural difficulty factors with a probe that asks each model to point at the pivot and outline the part separates grounding failures from structural and execution failures.

\paragraph{Recommendation.} We recommend, in order: (1) retain the focused benchmark and lead with direction (f), piloting the judge-blindness experiment on 50 clips by 5 December; (2) add the image-to-video comparison as a gated secondary experiment, after fixing a noise-floor protocol by 10 December, and drop it if it is not complete by 24 December; (3) reduce reinforcement learning to evidence that the reward resists exploitation and that grader-filtered best-of-$n$ sampling improves pass rates; (4) defer the cross-modal family to a subsequent paper, retaining only the single-frame edit control; and (5) adopt the side-effect vocabulary without the safety framing, citing action-level benchmarks as the policy counterpart to our capability measure. Regardless of direction, the paper must compare directly with MotionGraphicsBench in its introduction and must report human agreement for the grader, because MotionGraphicsBench reports 93.3\%.

\paragraph{Conditions that would change this recommendation.} If the pilot shows that frontier judges detect small side effects reliably, the judge-blindness experiment should leave the headline and direction (d) should move up. If another raster-input animation benchmark appears before January, differentiation would rest on grader validation and failure attribution alone.

\section{Trends Relevant to Timing}
Three trends, assessed from release notes, benchmark announcements and secondary coverage, bear on the timing of the paper. First, code-generated animation became a common demonstration and product category in late 2026, with five benchmarks released between September and November. Second, side effects of agents became a distinct evaluation axis with a growing number of benchmarks. Third, verifiers rather than benchmarks increasingly limit both evaluation and reinforcement learning, which raises the value of a validated, target-free grader. We judge all three signals to be strong; the first two also indicate that the space is becoming crowded, which argues for submitting in January rather than waiting. Details and sources are in \texttt{docs/research/trends.md}.

\bibliographystyle{plainnat}
\bibliography{refs}
\end{document}
